\begin{proof}[Proof of \cref{obj:thm:sequential-minimax}]
Throughout, nonnegative risks are interpreted in \([0,\infty]\). Set
\[
c=(-1,1)^\top,\qquad
J_0=J^*(\theta_0,p,\varepsilon),\qquad
V_0=V^*(\theta_0,p,\varepsilon),
\qquad
\iota(a)=\max\{a,0\}\in[0,\infty]\quad(a\in\mathbb R).
\]
For any procedure sequence \(\mathcal P\), write
\[
\operatorname{Risk}_{n,H}(\mathcal P)
=
\sup_{h\in\mathcal H_{n,H}(\theta_0)}
\mathbb E_{\theta_{n,h}}^{\mathcal P}
\!\left[(U_{n,h}^{\mathcal P})^2\right],
\qquad
\operatorname{Risk}_{\infty}(\mathcal P)
=
\sup_{H>0}\liminf_{n\to\infty}
\operatorname{Risk}_{n,H}(\mathcal P).
\]

\begin{enumerate}
\item
Fix an arbitrary sequentially private procedure sequence \(\mathcal P\), including its output spaces and estimator. Take the oracle coordinate \(t^*\) supplied by \cref{obj:thm:finite-oracle}, as used in \cref{obj:def:sequential-vt-handle}, and define
\[
v^*=v(t^*)=(t^*,t^*+1)^\top,
\qquad
G(\theta',t)
=
\max_{\alpha\in\mathcal A_\varepsilon}
F_{\theta'}(\alpha,t)
\quad
\bigl(\theta'\in(0,1)^2,\ t\in\mathbb R\bigr).
\]
Then
\[
c^\top v^*=1,
\qquad
G(\theta_0,t^*)=J_0.
\]
In particular, \(v^*\ne0\). For \(R>0\), put
\[
H(R)=R\|v^*\|_2>0,
\qquad
\vartheta_n(a)=\theta_0+\frac{a v^*}{\sqrt n}
\quad(n\ge1,\ a\in\mathbb R).
\]
For every fixed \(R\), the entire path \(\vartheta_n([-R,R])\) is interior for sufficiently large \(n\), because
\[
\sup_{|a|\le R}\|\vartheta_n(a)-\theta_0\|_2
\le \frac{R\|v^*\|_2}{\sqrt n}\longrightarrow0.
\]

Use the prior family in \cref{obj:def:sequential-vt-handle} at half the path radius:
\[
\overline w_R(a)=w_{R/2}(a)
=
\begin{cases}
\displaystyle
\frac{15}{8R}\left(1-\frac{4a^2}{R^2}\right)^2,
& |a|<R/2,\\[4pt]
0,& |a|\ge R/2.
\end{cases}
\]
Its derivative is
\[
\overline w_R'(a)
=
\begin{cases}
\displaystyle
-\frac{30a}{R^3}\left(1-\frac{4a^2}{R^2}\right),
& |a|<R/2,\\[4pt]
0,& |a|\ge R/2.
\end{cases}
\]
Thus \(\overline w_R\) is continuously differentiable, is nonnegative, and vanishes at both endpoints of \([-R,R]\). Direct integration gives
\[
\int_{-R}^{R}\overline w_R(a)\,da=1,
\qquad
\int_{-R}^{R}
\frac{\overline w_R'(a)^2}{\overline w_R(a)}\,da
=
\frac{480}{R^5}\int_{-R/2}^{R/2}a^2\,da
=
\frac{40}{R^2},
\]
where the quotient is zero wherever \(\overline w_R=0\).

For sufficiently large \(n\), define the prior-averaged information envelope
\[
\mathcal J_{n,R}
=
\int_{-R}^{R}
\overline w_R(a)G(\vartheta_n(a),t^*)\,da.
\]
The feasible weight set is compact: its row-normalization equations and nonnegative coordinates imply
\[
0\le\alpha_S\le1
\quad(S\in\mathcal S,\ \alpha\in\mathcal A_\varepsilon).
\]
Moreover, for interior \(\theta'\),
\[
1\le h_S(\theta')\le e^\varepsilon.
\]
Consequently, \(F_{\theta'}(\alpha,t^*)\) is jointly continuous in \((\theta',\alpha)\). Compactness gives continuity of its maximum, with
\[
|G(\theta',t^*)-G(\theta_0,t^*)|
\le
\max_{\alpha\in\mathcal A_\varepsilon}
|F_{\theta'}(\alpha,t^*)-F_{\theta_0}(\alpha,t^*)|.
\]
It follows that
\[
\sup_{|a|\le R}
|G(\vartheta_n(a),t^*)-J_0|
\longrightarrow0,
\qquad
\mathcal J_{n,R}\longrightarrow J_0.
\]
Equivalently, one may integrate the pointwise convergence using the eventual dominating function
\[
\bigl|\overline w_R(a)G(\vartheta_n(a),t^*)\bigr|
\le (|J_0|+1)\overline w_R(a).
\]
% lean: CausalSmith.Stat.LdpAteEfficiencySurface.sequentialVTPrior_facts
% lean: CausalSmith.Stat.LdpAteEfficiencySurface.tendsto_sequentialPriorInformationEnvelope

\item
We next obtain the transcript information bound. The product record law in \cref{obj:def:binary-trial-model}, with the category ordering in \cref{obj:synth_10}, has probabilities
\[
\begin{aligned}
\pi_{\theta',0}&=(1-p)(1-\theta'_0),&
\pi_{\theta',1}&=(1-p)\theta'_0,\\
\pi_{\theta',2}&=p(1-\theta'_1),&
\pi_{\theta',3}&=p\theta'_1
\end{aligned}
\qquad(\theta'\in(0,1)^2).
\]
Their directional derivatives along \(v^*\) are
\[
\dot\pi_j^*=(v^*)^\top\nabla_{\theta'}\pi_{\theta',j}
\quad(j=0,1,2,3),
\qquad
\sum_{j=0}^3\dot\pi_j^*=0.
\]

Fix \(n\ge1\). For \(0\le k\le n\), let
\[
\mathsf H_{n,k}=\prod_{i=1}^{k}\mathcal Z_{n,i}.
\]
For each input-category string
\(\mathbf j=(j_1,\ldots,j_k)\in\{0,1,2,3\}^k\), let
\(\Lambda_{n,k,\mathbf j}\) be the probability measure on \(\mathsf H_{n,k}\) obtained by successively applying the first \(k\) release kernels to that fixed input string. Define
\[
\nu_{n,k}
=
\sum_{\mathbf j\in\{0,1,2,3\}^k}\Lambda_{n,k,\mathbf j},
\qquad
\lambda_{n,k,\mathbf j}
=
\frac{d\Lambda_{n,k,\mathbf j}}{d\nu_{n,k}}.
\]
At \(k=0\), the history space is a singleton and its law is the unit point mass. Factorization of the transcript laws in \cref{obj:def:sequential-channels} yields the prefix density
\[
f_{n,k}^{\theta'}(\mathsf y)
=
\sum_{\mathbf j\in\{0,1,2,3\}^k}
\left(\prod_{i=1}^{k}\pi_{\theta',j_i}\right)
\lambda_{n,k,\mathbf j}(\mathsf y).
\]
Its directional derivative is the finite sum
\[
\dot f_{n,k}^{\theta'}(\mathsf y)
=
\sum_{\mathbf j\in\{0,1,2,3\}^k}
\left[
\sum_{i=1}^{k}
\dot\pi_{j_i}^*
\prod_{\substack{1\le i'\le k\\i'\ne i}}
\pi_{\theta',j_{i'}}
\right]
\lambda_{n,k,\mathbf j}(\mathsf y).
\]
Let \(\mathbb P_{n,k}^{\theta'}\) denote this prefix law, and define its score by
\[
s_{n,k}^{\theta'}(\mathsf y)
=
\begin{cases}
\dot f_{n,k}^{\theta'}(\mathsf y)/f_{n,k}^{\theta'}(\mathsf y),
&f_{n,k}^{\theta'}(\mathsf y)>0,\\
0,&f_{n,k}^{\theta'}(\mathsf y)=0.
\end{cases}
\]
All record probabilities are positive. Hence these scores are square integrable; indeed,
\[
|s_{n,k}^{\theta'}|
\le
k\max_{0\le j\le3}
\frac{|\dot\pi_j^*|}{\pi_{\theta',j}}
\quad
\mathbb P_{n,k}^{\theta'}\text{-almost everywhere}.
\]

For \(k<n\), freeze the prefix law at \(\theta'\) and form the joint channel
\[
\overline Q_{n,k}^{\theta'}(d\mathsf y,dz\mid x)
=
\mathbb P_{n,k}^{\theta'}(d\mathsf y)\,
Q_{k+1}^{(n)}(dz\mid x,\mathsf y).
\]
For a measurable joint event \(A\), write
\[
A_{\mathsf y}=\{z:(\mathsf y,z)\in A\}.
\]
The all-history privacy inequalities give
\[
\overline Q_{n,k}^{\theta'}(A\mid x)
=
\int Q_{k+1}^{(n)}(A_{\mathsf y}\mid x,\mathsf y)
\,\mathbb P_{n,k}^{\theta'}(d\mathsf y)
\le
e^\varepsilon\overline Q_{n,k}^{\theta'}(A\mid x').
\]
Thus this is a stationary private channel on the joint history--release space. Define its output law and directional score by
\[
\rho_{n,k}^{\theta'}
=
\sum_{j=0}^3
\pi_{\theta',j}\overline Q_{n,k}^{\theta'}(\,\cdot\mid x_j),
\qquad
\Delta_{n,k}^{\theta'}
=
\frac{d\left(\sum_{j=0}^3
\dot\pi_j^*\overline Q_{n,k}^{\theta'}(\,\cdot\mid x_j)\right)}
{d\rho_{n,k}^{\theta'}}.
\]
Choose a measurable version of the latter derivative, with value zero on a null set.

Differentiating the finite-mixture expression for the successor prefix law gives, almost everywhere under \(\rho_{n,k}^{\theta'}\),
\[
s_{n,k+1}^{\theta'}(\mathsf y,z)
=
s_{n,k}^{\theta'}(\mathsf y)
+
\Delta_{n,k}^{\theta'}(\mathsf y,z).
\]
Here the first summand differentiates the prefix law and the second differentiates the current-record mixture. For every measurable prefix event \(B\),
\[
\int_{B\times\mathcal Z_{n,k+1}}
\Delta_{n,k}^{\theta'}\,d\rho_{n,k}^{\theta'}
=
\sum_{j=0}^3\dot\pi_j^*\,
\mathbb P_{n,k}^{\theta'}(B)
=
0.
\]
Therefore the current score has conditional mean zero given the prefix, and
\[
\int
s_{n,k}^{\theta'}(\mathsf y)\Delta_{n,k}^{\theta'}(\mathsf y,z)
\,d\rho_{n,k}^{\theta'}(\mathsf y,z)
=
0.
\]

Define the prefix score energy
\[
\mathcal I_{n,k}(\theta')
=
\int (s_{n,k}^{\theta'})^2\,d\mathbb P_{n,k}^{\theta'}.
\]
Expanding the square in the score identity gives
\[
\mathcal I_{n,0}(\theta')=0,
\qquad
\mathcal I_{n,k+1}(\theta')
=
\mathcal I_{n,k}(\theta')
+
\int(\Delta_{n,k}^{\theta'})^2\,d\rho_{n,k}^{\theta'}.
\]
Apply \cref{obj:lem:staircase-refinement} to the joint channel. It supplies feasible weights
\(\beta_{n,k}^{\theta'}\in\mathcal A_\varepsilon\) with
\[
\begin{aligned}
\int(\Delta_{n,k}^{\theta'})^2\,d\rho_{n,k}^{\theta'}
&=(v^*)^\top
I_{\theta'}(\overline Q_{n,k}^{\theta'})v^*\\
&\le F_{\theta'}(\beta_{n,k}^{\theta'},t^*)
\le G(\theta',t^*).
\end{aligned}
\]
Induction on the prefix length now yields
\[
\mathcal I_{n,k}(\theta')\le kG(\theta',t^*)
\quad(0\le k\le n).
\]
For the local scalar path \(a\mapsto\vartheta_n(a)\), define
\[
s_n^{\mathrm{loc}}(a,z)
=
\frac{1}{\sqrt n}s_{n,n}^{\vartheta_n(a)}(z),
\qquad
I_n^{\mathrm{loc}}(a)
=
\int\bigl(s_n^{\mathrm{loc}}(a,z)\bigr)^2
\,\mathbb P_{n,n}^{\vartheta_n(a)}(dz).
\]
For sufficiently large \(n\), uniformly over \(|a|\le R\),
\[
0\le I_n^{\mathrm{loc}}(a)
=
\frac{\mathcal I_{n,n}(\vartheta_n(a))}{n}
\le G(\vartheta_n(a),t^*).
\]
Multiplying by the nonnegative prior density and integrating gives
\[
0\le
\int_{-R}^{R}\overline w_R(a)I_n^{\mathrm{loc}}(a)\,da
\le\mathcal J_{n,R}.
\]
The finite-mixture domination and the displayed score bounds also give the integrability needed for this calculation on arbitrary measurable output spaces.
% lean: CausalSmith.Stat.LdpAteEfficiencySurface.prefixScoreEnergy_succ
% lean: CausalSmith.Stat.LdpAteEfficiencySurface.prefixScoreEnergy_le_mul_upperEnvelope
% lean: CausalSmith.Stat.LdpAteEfficiencySurface.scaledSelected_fisher_integrable_and_control

\item
Fix \(R>0\) and take \(n\) sufficiently large that the preceding path is interior. Define the prior--transcript probability measure
\[
\mathbb B_{n,R}(da,dz)
=
\overline w_R(a)\,da\,
\mathbb P_{n,n}^{\vartheta_n(a)}(dz)
\quad(a\in[-R,R]),
\]
the estimation error and its integrated square by
\[
e_n(a,z)=\widehat\tau_n^{\mathcal P}(z)-\tau(\vartheta_n(a)),
\qquad
\mathcal E_{n,R}=\int e_n(a,z)^2\,\mathbb B_{n,R}(da,dz),
\]
and the joint score by
\[
\mathscr S_{n,R}(a,z)
=
s_n^{\mathrm{loc}}(a,z)
+
\begin{cases}
\overline w_R'(a)/\overline w_R(a),&\overline w_R(a)>0,\\
0,&\overline w_R(a)=0.
\end{cases}
\]

First suppose \(\mathcal E_{n,R}<\infty\). The likelihood score has mean zero at each path parameter. Hence its cross term with the prior score integrates to zero, and
\[
\int\mathscr S_{n,R}^2\,d\mathbb B_{n,R}
=
\int_{-R}^{R}\overline w_R(a)I_n^{\mathrm{loc}}(a)\,da
+
\frac{40}{R^2}.
\]
This quantity is positive and is at most \(\mathcal J_{n,R}+40/R^2\).

The target derivative along the path is
\[
\frac{d}{da}\tau(\vartheta_n(a))
=
\frac{c^\top v^*}{\sqrt n}
=
\frac1{\sqrt n}.
\]
Integration by parts therefore gives
\[
\int e_n\mathscr S_{n,R}\,d\mathbb B_{n,R}
=
\frac1{\sqrt n}.
\]
Indeed, the derivative of the product of the error and the joint density is the error times the joint-density derivative minus \(1/\sqrt n\) times that density. Its boundary integral vanishes because \(\overline w_R(\pm R)=0\). Finite integrated squared error and finite joint score energy justify the integrals by Cauchy--Schwarz. Applying the same inequality to the displayed identity yields
\[
\begin{aligned}
n\mathcal E_{n,R}
&\ge
\frac{1}{
\displaystyle
\int_{-R}^{R}\overline w_R(a)I_n^{\mathrm{loc}}(a)\,da
+40/R^2}\\
&\ge
\frac1{\mathcal J_{n,R}+40/R^2}.
\end{aligned}
\]

For every \(|a|\le R\),
\[
av^*\in\mathcal H_{n,H(R)}(\theta_0),
\qquad
\sqrt n\,e_n(a,z)=U_{n,av^*}^{\mathcal P}(z).
\]
Thus the normalized prior average is bounded by the local worst risk:
\[
\begin{aligned}
\iota(n\mathcal E_{n,R})
&=
\int_{-R}^{R}\overline w_R(a)
\mathbb E_{\vartheta_n(a)}^{\mathcal P}
\!\left[(U_{n,av^*}^{\mathcal P})^2\right]\,da\\
&\le \operatorname{Risk}_{n,H(R)}(\mathcal P).
\end{aligned}
\]
If instead \(\mathcal E_{n,R}=\infty\), the same nonnegative-integral identity makes this prior average infinite. Its domination by the local worst risk then gives
\(\operatorname{Risk}_{n,H(R)}(\mathcal P)=\infty\).
Both cases establish the eventual bound
\[
\iota\!\left(\frac1{\mathcal J_{n,R}+40/R^2}\right)
\le
\operatorname{Risk}_{n,H(R)}(\mathcal P).
\]
% lean: CausalSmith.Stat.LdpAteEfficiencySurface.ofReal_vanTreesEnvelope_le_localWorstRisk_of_fisher_control_allErrors

\item
We record positivity of the optimized information before passing to reciprocal limits. Define
\[
a_{\mathrm{ctl}}
=
\frac{(e^\varepsilon-1)^2(1-p)^2}
{e^\varepsilon(e^\varepsilon+3)},
\qquad
a_{\mathrm{trt}}
=
\frac{(e^\varepsilon-1)^2p^2}
{e^\varepsilon(e^\varepsilon+3)},
\qquad
J_{\mathrm{low}}
=
\frac{a_{\mathrm{ctl}}a_{\mathrm{trt}}}
{a_{\mathrm{ctl}}+a_{\mathrm{trt}}}.
\]
These constants are positive. The singleton staircase design
\[
\alpha^{\mathrm{RR}}_S
=
\begin{cases}
(e^\varepsilon+3)^{-1},&S\text{ is a singleton},\\
0,&\text{otherwise}
\end{cases}
\]
belongs to \(\mathcal A_\varepsilon\): each channel row contains one singleton entry with multiplier \(e^\varepsilon\) and three with multiplier one.

Using the singleton patterns \(\{x_0\}\) and \(\{x_3\}\), the upper bound \(h_S(\theta')\le e^\varepsilon\), and nonnegativity of the remaining terms gives, for every interior \(\theta'\) and every \(t\in\mathbb R\),
\[
\begin{aligned}
F_{\theta'}(\alpha^{\mathrm{RR}},t)
&\ge
a_{\mathrm{ctl}}t^2+a_{\mathrm{trt}}(t+1)^2\\
&=
(a_{\mathrm{ctl}}+a_{\mathrm{trt}})
\left(t+\frac{a_{\mathrm{trt}}}
{a_{\mathrm{ctl}}+a_{\mathrm{trt}}}\right)^2
+J_{\mathrm{low}}
\ge J_{\mathrm{low}}.
\end{aligned}
\]
The optimization in \cref{obj:def:staircase-saddle} consequently implies
\[
J^*(\theta',p,\varepsilon)\ge J_{\mathrm{low}}>0,
\qquad
J_0>0.
\]

For each fixed \(R>0\), the convergence of \(\mathcal J_{n,R}\) and the eventual lower bound therefore imply
\[
\iota\!\left(\frac1{J_0+40/R^2}\right)
\le
\liminf_{n\to\infty}
\operatorname{Risk}_{n,H(R)}(\mathcal P)
\le
\operatorname{Risk}_{\infty}(\mathcal P).
\]
To remove the prior-information penalty, take
\[
R_{\mathrm{rad},\ell}=\ell+1\quad(\ell\in\mathbb N).
\]
Then
\[
\frac{40}{R_{\mathrm{rad},\ell}^2}\longrightarrow0,
\qquad
\frac1{J_0+40/R_{\mathrm{rad},\ell}^2}\longrightarrow\frac1{J_0}=V_0.
\]
Passing to this limit proves
\[
\iota(V_0)\le\operatorname{Risk}_{\infty}(\mathcal P).
\]
Since \(\mathcal P\) and its measurable output spaces were arbitrary, this establishes the universal lower bound.
% lean: CausalSmith.Stat.LdpAteEfficiencySurface.inv_information_le_localAsymptoticRisk_of_priorRadius_family
% lean: CausalSmith.Stat.LdpAteEfficiencySurface.sequential_localAsymptoticRisk_lower_bound

\item
Choose the schedule
\[
m_n=\lfloor\sqrt n\rfloor\quad(n\in\mathbb N).
\]
For every integer \(b\ge0\),
\[
n\ge b^2\quad\Longrightarrow\quad m_n\ge b,
\]
so \(m_n\to\infty\). For \(n\ge1\),
\[
0\le\frac{m_n}{n}
\le\frac{\sqrt n}{n}
=\frac1{\sqrt n}\longrightarrow0.
\]
Also,
\[
1\le m_n<n\qquad(n\ge2).
\]
Thus this schedule satisfies
\cref{obj:ass:pilot-diverges,obj:ass:pilot-sublinear}.
% lean: CausalSmith.Stat.LdpAteEfficiencySurface.natSqrt_pilotRates

\item
Let
\[
\mathcal P^*
=
\text{the procedure in \cref{obj:def:pilot-estimator}
with this schedule and the fixed strong-saddle selector}.
\]
Apply \cref{obj:thm:pilot-attainment} at \(\theta_0\), with \(\gamma=1/2\). It supplies sequential privacy at every sample size, the common local limit
\[
L_{\mathcal P^*}=N(0,V_0),
\qquad
\int_{\mathbb R}u^2\,L_{\mathcal P^*}(du)=V_0<\infty,
\]
and, for every \(H>0\),
\[
\lim_{\substack{M\to\infty\\M\in\mathbb N}}
\limsup_{n\to\infty}
\sup_{h\in\mathcal H_{n,H}(\theta_0)}
\mathbb E_{\theta_{n,h}}^{\mathcal P^*}
\!\left[
(U_{n,h}^{\mathcal P^*})^2
\mathbf1\{|U_{n,h}^{\mathcal P^*}|>M\}
\right]
=0.
\]
Monotonicity in the threshold extends this limit to real \(M\to\infty\). These are precisely the regularity requirements, so
\[
\mathcal P^*\in\mathfrak P^{\mathrm{reg}}_\varepsilon(\theta_0).
\]
The conventions in \cref{obj:def:pilot-estimator} specify the procedure at every sample size. The following risk calculations use \(n\ge2\), when both sample portions are nonempty.
% lean: CausalSmith.Stat.LdpAteEfficiencySurface.pilot_attainment

\item
We establish exact risk attainment through the pilot-averaged conditional variance. First obtain uniform bounds for the selected correction. Define
\[
C_{\mathrm{zero}}
=
\sum_{S\in\mathcal S}(g_S^\top v(0))^2,
\qquad
T_{\mathrm{up}}
=
\sqrt{\frac{C_{\mathrm{zero}}}{a_{\mathrm{ctl}}}},
\]
\[
C_{\mathrm{grad}}
=
\sum_{S\in\mathcal S}(|g_{S,0}|+|g_{S,1}|),
\qquad
C_\phi
=
\frac{C_{\mathrm{grad}}(T_{\mathrm{up}}+1)}
{J_{\mathrm{low}}}.
\]
For an interior selector argument \(\eta\), set
\[
v_\eta=v(t_\eta),
\qquad
\mathcal I_\eta=I_\eta(\alpha_\eta),
\qquad
\phi_\eta(S)=
\frac{g_S^\top v_\eta}{J_\eta h_S(\eta)}
\quad(S\in\mathcal S).
\]
The strong-saddle conditions in \cref{obj:def:pilot-estimator} and the reference-design bound give
\[
J_{\mathrm{low}}
\le
a_{\mathrm{ctl}}t_\eta^2+a_{\mathrm{trt}}(t_\eta+1)^2
\le
F_\eta(\alpha^{\mathrm{RR}},t_\eta)
\le
J_\eta.
\]
On the other hand, minimizing in the profiling coordinate and using
\(\alpha_{\eta,S}\le1\) and \(h_S(\eta)\ge1\) gives
\[
J_\eta
=
F_\eta(\alpha_\eta,t_\eta)
\le F_\eta(\alpha_\eta,0)
\le C_{\mathrm{zero}}.
\]
It follows that
\[
|t_\eta|\le T_{\mathrm{up}},
\qquad
J_\eta\ge J_{\mathrm{low}}.
\]
Finally,
\[
|g_S^\top v_\eta|
\le
(|g_{S,0}|+|g_{S,1}|)(|t_\eta|+1)
\le
C_{\mathrm{grad}}(T_{\mathrm{up}}+1),
\]
and hence
\[
|\phi_\eta(S)|\le C_\phi
\quad(\eta\in(0,1)^2,\ S\in\mathcal S).
\]
% lean: CausalSmith.Stat.LdpAteEfficiencySurface.strongSelector_J_bounds
% lean: CausalSmith.Stat.LdpAteEfficiencySurface.strongSelector_t_bound
% lean: CausalSmith.Stat.LdpAteEfficiencySurface.strongSelector_phiTilde_bound

\item
For interior true and pilot parameters \(\theta'\) and \(\eta\), define
\[
\sigma^2(\theta',\eta)
=
\sum_{S\in\mathcal S}
\alpha_{\eta,S}h_S(\theta')\phi_\eta(S)^2
-
\bigl(\tau(\theta')-\tau(\eta)\bigr)^2.
\]
We derive the mean and second-moment identities entering this expression. Put
\[
\mathbf{1}_2=(1,1)^\top.
\]
Since \(v(t)=v(0)+t\mathbf{1}_2\), minimization of the selected quadratic objective gives
\[
0=
\left.\frac{d}{dt}F_\eta(\alpha_\eta,t)\right|_{t=t_\eta}
=
2\mathbf{1}_2^\top\mathcal I_\eta v_\eta.
\]
Thus \(\mathcal I_\eta v_\eta\) lies in the span of \(c\). Together with
\(c^\top v_\eta=1\) and \(v_\eta^\top\mathcal I_\eta v_\eta=J_\eta\), this yields
\[
\mathcal I_\eta v_\eta=J_\eta c.
\]

Channel normalization and the affine dependence of the pattern masses give
\[
\sum_S\alpha_{\eta,S}h_S(\theta')=1,
\qquad
\sum_S\alpha_{\eta,S}g_S=0,
\qquad
h_S(\theta')=h_S(\eta)+g_S^\top(\theta'-\eta).
\]
Therefore
\[
\begin{aligned}
\sum_S\alpha_{\eta,S}h_S(\theta')\phi_\eta(S)
&=
\frac1{J_\eta}
\left[
\sum_S\alpha_{\eta,S}g_S^\top v_\eta
+
(\theta'-\eta)^\top\mathcal I_\eta v_\eta
\right]\\
&=
c^\top(\theta'-\eta)
=
\tau(\theta')-\tau(\eta).
\end{aligned}
\]
At the diagonal,
\[
\begin{aligned}
\sum_S\alpha_{\eta,S}h_S(\eta)\phi_\eta(S)^2
&=
\frac1{J_\eta^2}
\sum_S
\alpha_{\eta,S}\frac{(g_S^\top v_\eta)^2}{h_S(\eta)}\\
&=\frac1{J_\eta}=V^*(\eta,p,\varepsilon).
\end{aligned}
\]
Consequently,
\[
\sigma^2(\theta',\eta)
=
\sum_S\alpha_{\eta,S}h_S(\theta')
\left[
\phi_\eta(S)-\bigl(\tau(\theta')-\tau(\eta)\bigr)
\right]^2,
\]
and
\[
0\le\sigma^2(\theta',\eta)\le C_\phi^2,
\qquad
\sigma^2(\eta,\eta)=V^*(\eta,p,\varepsilon).
\]
The upper bound follows by bounding the uncentered second moment by \(C_\phi^2\) and subtracting the squared mean.
% lean: CausalSmith.Stat.LdpAteEfficiencySurface.phiTilde_mean_identity
% lean: CausalSmith.Stat.LdpAteEfficiencySurface.phiTilde_secondMoment_identity
% lean: CausalSmith.Stat.LdpAteEfficiencySurface.adaptiveScoreVariance_nonneg
% lean: CausalSmith.Stat.LdpAteEfficiencySurface.adaptiveScoreVariance_diag

\item
For \(n\ge2\), define the main-sample size and pilot-prefix probabilities by
\[
N_n=n-m_n>0,
\]
\[
\omega_{n,\theta'}(\mathbf r)
=
\prod_{i=1}^{m_n}
\left[
\sum_{j=0}^3
\pi_{\theta',j}Q_{\mathrm{RR}}(r_i\mid x_j)
\right],
\qquad
\mathbf r\in\{0,1,2,3\}^{m_n},
\quad \theta'\in(0,1)^2.
\]
They satisfy
\[
\omega_{n,\theta'}(\mathbf r)\ge0,
\qquad
\sum_{\mathbf r\in\{0,1,2,3\}^{m_n}}
\omega_{n,\theta'}(\mathbf r)=1.
\]
Let \(\widetilde\theta(\mathbf r)\) be the clipped pilot estimate defined in \cref{obj:def:pilot-estimator}. For a main-output string
\(\mathbf s\in\mathcal S^{N_n}\), factorization gives
\[
\begin{aligned}
&\Pr_{\theta'}^{\mathcal P^*}
\!\left(H_{m_n}=\mathbf r,\,
(S_{m_n+1},\ldots,S_n)=\mathbf s\right)\\
&\qquad=
\omega_{n,\theta'}(\mathbf r)
\prod_{i=1}^{N_n}
\alpha_{\widetilde\theta(\mathbf r),s_i}
h_{s_i}(\theta').
\end{aligned}
\]
Transcripts outside these pilot and main-output formats have probability zero.

Given the pilot prefix, the main corrections are independent with mean
\(\tau(\theta')-\tau(\widetilde\theta(\mathbf r))\) and variance
\(\sigma^2(\theta',\widetilde\theta(\mathbf r))\). The estimator formula therefore gives
\[
\mathbb E_{\theta'}^{\mathcal P^*}
\!\left[
\widehat\tau^*-\tau(\theta')
\mid H_{m_n}=\mathbf r
\right]=0.
\]
Expanding the square of the centered sum, the off-diagonal terms vanish by independence and zero means, so
\[
\mathbb E_{\theta'}^{\mathcal P^*}
\!\left[
n(\widehat\tau^*-\tau(\theta'))^2
\mid H_{m_n}=\mathbf r
\right]
=
\frac n{N_n}\,
\sigma^2(\theta',\widetilde\theta(\mathbf r)).
\]
Define the prefix-averaged variance
\[
\overline\sigma_n^2(\theta')
=
\sum_{\mathbf r\in\{0,1,2,3\}^{m_n}}
\omega_{n,\theta'}(\mathbf r)
\sigma^2(\theta',\widetilde\theta(\mathbf r)).
\]
Summing the conditional identity over the finite prefix space gives, for every admissible \(h\),
\[
\mathbb E_{\theta_{n,h}}^{\mathcal P^*}
\!\left[(U_{n,h}^{\mathcal P^*})^2\right]
=
\frac n{N_n}\,
\overline\sigma_n^2(\theta_{n,h}),
\qquad
0\le\overline\sigma_n^2(\theta_{n,h})\le C_\phi^2.
\]
% lean: CausalSmith.Stat.LdpAteEfficiencySurface.pilotEstimator_localRisk_eq_scaled_prefixVarianceAverage
% lean: CausalSmith.Stat.LdpAteEfficiencySurface.pilotPrefixVarianceAverage_bounds

\item
We next show continuity of the variance at \((\theta_0,\theta_0)\). The uniform selected-coordinate bound supplies the compact interval
\[
K_{\mathrm{prof}}=[-T_{\mathrm{up}},T_{\mathrm{up}}].
\]
For every interior \(\eta\), the selected saddle lies in
\(\mathcal A_\varepsilon\times K_{\mathrm{prof}}\). Its saddle inequalities imply
\[
J^*(\eta,p,\varepsilon)
=
\max_{\alpha\in\mathcal A_\varepsilon}
\min_{t\in K_{\mathrm{prof}}}F_\eta(\alpha,t).
\]
Indeed, evaluating any feasible design at \(t_\eta\) bounds its restricted minimum above by \(J_\eta\), while the selected design has restricted minimum \(J_\eta\).

Joint continuity on the fixed compact weight--coordinate set gives
\[
\begin{aligned}
|J^*(\eta,p,\varepsilon)-J_0|
&\le
\sup_{\substack{\alpha\in\mathcal A_\varepsilon\\
t\in K_{\mathrm{prof}}}}
|F_\eta(\alpha,t)-F_{\theta_0}(\alpha,t)|
\longrightarrow0
\quad(\eta\to\theta_0).
\end{aligned}
\]
Since \(J_0>0\), reciprocation yields
\[
V^*(\eta,p,\varepsilon)\longrightarrow V_0.
\]

Subtracting the diagonal second-moment identity from the definition of \(\sigma^2\) gives the exact expression
\[
\begin{aligned}
\sigma^2(\theta',\eta)-V^*(\eta,p,\varepsilon)
&=
\sum_S
\alpha_{\eta,S}\bigl(h_S(\theta')-h_S(\eta)\bigr)
\phi_\eta(S)^2\\
&\qquad-
\bigl(\tau(\theta')-\tau(\eta)\bigr)^2.
\end{aligned}
\]
Using \(0\le\alpha_{\eta,S}\le1\) and \(|\phi_\eta(S)|\le C_\phi\), we obtain
\[
\begin{aligned}
|\sigma^2(\theta',\eta)-V^*(\eta,p,\varepsilon)|
&\le
C_\phi^2\sum_S|h_S(\theta')-h_S(\eta)|\\
&\qquad+
\bigl(\tau(\theta')-\tau(\eta)\bigr)^2.
\end{aligned}
\]
The right-hand side tends to zero as
\((\theta',\eta)\to(\theta_0,\theta_0)\). Hence
\[
\sigma^2(\theta',\eta)\longrightarrow V_0
\quad\text{as }(\theta',\eta)\to(\theta_0,\theta_0).
\]
% lean: CausalSmith.Stat.LdpAteEfficiencySurface.continuousAt_Jstar
% lean: CausalSmith.Stat.LdpAteEfficiencySurface.continuousAt_Vstar
% lean: CausalSmith.Stat.LdpAteEfficiencySurface.continuousAt_adaptiveScoreVariance_diag

\item
Fix \(H>0\) and an error tolerance \(\delta>0\). By the preceding continuity, choose \(r_{\mathrm{cont}}>0\) such that
\[
\|\theta'-\theta_0\|_\infty<r_{\mathrm{cont}},
\quad
\|\eta-\theta_0\|_\infty<r_{\mathrm{cont}}
\quad\Longrightarrow\quad
|\sigma^2(\theta',\eta)-V_0|<\delta/2
\]
for interior \(\theta'\) and \(\eta\). Define
\[
d_{\mathrm{int}}
=
\min\{\mu_0,1-\mu_0,\mu_1,1-\mu_1\}>0,
\qquad
r_{\mathrm{pil}}
=
\min\{r_{\mathrm{cont}}/3,d_{\mathrm{int}}/8\}>0,
\]
where \(\theta_0=(\mu_0,\mu_1)^\top\).

Uniformly over the bounded local set,
\[
\|\theta_{n,h}-\theta_0\|_\infty
\le\frac H{\sqrt n}.
\]
Thus, eventually, this distance is less than \(r_{\mathrm{pil}}\), and
\[
\delta_n=(m_n+2)^{-1}<d_{\mathrm{int}}/4.
\]
For these \(n\), every admissible local parameter satisfies
\[
\delta_n+r_{\mathrm{pil}}
\le
\min\{
(\theta_{n,h})_0,1-(\theta_{n,h})_0,
(\theta_{n,h})_1,1-(\theta_{n,h})_1
\}.
\]

For an interior \(\theta'\), define the bad-prefix mass
\[
\begin{aligned}
B_n(\theta';r)
&=
\sum_{\mathbf r\in\{0,1,2,3\}^{m_n}}
\omega_{n,\theta'}(\mathbf r)\,
\mathbf1\{
\|\widetilde\theta(\mathbf r)-\theta'\|_\infty\ge r
\}\\
&=
\Pr_{\theta'}^{\mathcal P^*}
\{\|\widetilde\theta-\theta'\|_\infty\ge r\},
\qquad r>0.
\end{aligned}
\]
The randomized-response probabilities in \cref{obj:synth_10} give
\[
\mathbb E_{\theta'}[\widehat u_1]
=
\frac{1+(e^\varepsilon-1)(1-p)\theta'_0}
{e^\varepsilon+3},
\qquad
\mathbb E_{\theta'}[\widehat u_3]
=
\frac{1+(e^\varepsilon-1)p\theta'_1}
{e^\varepsilon+3},
\]
and independence of the pilot releases gives
\[
\operatorname{Var}_{\theta'}(\widehat u_1)\le\frac1{m_n},
\qquad
\operatorname{Var}_{\theta'}(\widehat u_3)\le\frac1{m_n}.
\]
Define the positive frequency thresholds
\[
b_{\mathrm{ctl}}(r)
=
\frac{r(1-p)(e^\varepsilon-1)}{e^\varepsilon+3},
\qquad
b_{\mathrm{trt}}(r)
=
\frac{rp(e^\varepsilon-1)}{e^\varepsilon+3}.
\]
If both empirical-frequency deviations are smaller than their respective thresholds, the inverted arm estimates are within \(r\) of the true coordinates. Under the displayed interior-margin condition, clipping preserves these strict inequalities. Therefore Chebyshev's inequality and the union bound give
\[
B_n(\theta';r)
\le
\frac1{m_n b_{\mathrm{ctl}}(r)^2}
+
\frac1{m_n b_{\mathrm{trt}}(r)^2}.
\]
For fixed \(r>0\), this upper bound tends to zero as \(m_n\to\infty\).

On a good prefix at \(\theta'=\theta_{n,h}\), the pilot estimate satisfies
\[
\|\widetilde\theta(\mathbf r)-\theta_0\|_\infty
<
2r_{\mathrm{pil}}<r_{\mathrm{cont}},
\]
so its conditional variance differs from \(V_0\) by less than \(\delta/2\). On every prefix, the bound \(0\le\sigma^2\le C_\phi^2\) gives
\[
|\sigma^2(\theta_{n,h},\widetilde\theta(\mathbf r))-V_0|
\le C_\phi^2+|V_0|.
\]
Splitting the finite weighted average into good and bad prefixes consequently yields
\[
\begin{aligned}
|\overline\sigma_n^2(\theta_{n,h})-V_0|
&\le
\frac{\delta}{2}
+
(C_\phi^2+|V_0|)B_n(\theta_{n,h};r_{\mathrm{pil}})\\
&\le
\frac{\delta}{2}
+
(C_\phi^2+|V_0|)
\left[
\frac1{m_n b_{\mathrm{ctl}}(r_{\mathrm{pil}})^2}
+
\frac1{m_n b_{\mathrm{trt}}(r_{\mathrm{pil}})^2}
\right].
\end{aligned}
\]
The last term is eventually less than \(\delta/2\), uniformly over admissible \(h\). Hence
\[
\sup_{h\in\mathcal H_{n,H}(\theta_0)}
|\overline\sigma_n^2(\theta_{n,h})-V_0|
\longrightarrow0.
\]
% lean: CausalSmith.Stat.LdpAteEfficiencySurface.pilotPrefixBadMass_le
% lean: CausalSmith.Stat.LdpAteEfficiencySurface.pilotPrefixVarianceAverage_error_le
% lean: CausalSmith.Stat.LdpAteEfficiencySurface.eventually_uniform_pilotPrefixVarianceAverage

\item
The sublinear schedule gives
\[
\frac n{N_n}
=
\frac1{1-m_n/n}\longrightarrow1.
\]
Using the bound on the averaged variance,
\[
\begin{aligned}
\left|
\frac n{N_n}\overline\sigma_n^2(\theta_{n,h})-V_0
\right|
&\le
\left|\frac n{N_n}-1\right|C_\phi^2
+
|\overline\sigma_n^2(\theta_{n,h})-V_0|.
\end{aligned}
\]
Thus the exact risk identity implies
\[
\sup_{h\in\mathcal H_{n,H}(\theta_0)}
\left|
\mathbb E_{\theta_{n,h}}^{\mathcal P^*}
[(U_{n,h}^{\mathcal P^*})^2]-V_0
\right|
\longrightarrow0.
\]

In particular, for every \(\delta>0\), eventually
\[
\iota(V_0-\delta)
\le
\operatorname{Risk}_{n,H}(\mathcal P^*)
\le
\iota(V_0+\delta).
\]
For the lower inequality, use \(h=0\), which belongs to
\(\mathcal H_{n,H}(\theta_0)\); the upper inequality holds for every admissible perturbation and hence for their supremum. These bounds give convergence in the extended nonnegative reals:
\[
\operatorname{Risk}_{n,H}(\mathcal P^*)\longrightarrow\iota(V_0).
\]
The upper bounds also ensure that the risks are eventually finite. Consequently,
\[
\liminf_{n\to\infty}\operatorname{Risk}_{n,H}(\mathcal P^*)
=
\iota(V_0)
\qquad(H>0).
\]
Taking the supremum over positive radii, a nonempty family whose values all equal \(\iota(V_0)\), gives
\[
\operatorname{Risk}_{\infty}(\mathcal P^*)=\iota(V_0).
\]
This is the required attaining witness for the fixed selector.
% lean: CausalSmith.Stat.LdpAteEfficiencySurface.eventually_uniform_scaledPrefixVarianceAverage
% lean: CausalSmith.Stat.LdpAteEfficiencySurface.eventually_localWorstRisk_sandwich
% lean: CausalSmith.Stat.LdpAteEfficiencySurface.pilot_localWorstRisk_tendsto
% lean: CausalSmith.Stat.LdpAteEfficiencySurface.pilot_localAsymptoticRisk_eq

\item
Finally, the universal lower bound applies to every member of
\(\mathfrak P^{\mathrm{reg}}_\varepsilon(\theta_0)\), so
\[
\iota(V_0)
\le
\inf_{\mathcal P\in\mathfrak P^{\mathrm{reg}}_\varepsilon(\theta_0)}
\operatorname{Risk}_{\infty}(\mathcal P).
\]
The regular procedure \(\mathcal P^*\) constructed above belongs to this class and has risk \(\iota(V_0)\). Therefore
\[
\inf_{\mathcal P\in\mathfrak P^{\mathrm{reg}}_\varepsilon(\theta_0)}
\operatorname{Risk}_{\infty}(\mathcal P)
\le
\operatorname{Risk}_{\infty}(\mathcal P^*)
=
\iota(V_0).
\]
The two inequalities establish the claimed infimum equality in \([0,\infty]\).
% lean: CausalSmith.Stat.LdpAteEfficiencySurface.sequential_minimax
\end{enumerate}
\end{proof}
